\chapter{Zusammenfügen und Erweitern}
\label{b06guide:ch:extensions}

\section{Funktionen auf Teilbereichen zusammenfügen}

\begin{definition}[Kompatible Familie und zusammengefügte Funktion]
Seien $(A_i)_{i\in I}$ eine gegebene Familie von Mengen und
$f_i\colon A_i\to B$ gegebene Funktionen. Die Familie heißt
\emph{kompatibel}, wenn
\[
 f_i(x)=f_j(x)\qquad
 (i,j\in I,\ x\in A_i\cap A_j).
\]
Für eine kompatible Familie bezeichnet
\[
 f=\bigcup_{i\in I}f_i\colon\bigcup_{i\in I}A_i\to B
\]
die zusammengefügte Funktion mit $f(x)=f_i(x)$ für $x\in A_i$.
Die Vereinigung wird hier auf den Funktionsgraphen gebildet.
\end{definition}

\begin{theorem}[Injektives Zusammenfügen auf überlappenden Bereichen]
\label{b06guide:thm:gluing}
Sei $(f_i\colon A_i\to B)_{i\in I}$ eine gegebene kompatible
Familie. Ihre zusammengefügte Funktion $f=\bigcup_{i\in I}f_i$
ist genau dann injektiv, wenn beide folgenden Bedingungen gelten:
\begin{enumerate}
\item Jedes $f_i$ ist injektiv.
\item Für alle $i,j\in I$ gilt
\[
 f_i[A_i]\cap f_j[A_j]=f_i[A_i\cap A_j].
\]
\end{enumerate}
Wegen der Kompatibilität ist auf der rechten Seite auch
$f_j[A_i\cap A_j]$ zulässig.
\end{theorem}

\begin{theorem}[Injektives Zusammenfügen auf disjunkten Bereichen]
\label{b06guide:thm:disjointgluing}
Seien die Mengen $A_i$, $i\in I$, paarweise disjunkt, und
seien $f_i\colon A_i\rightarrowtail B$ gegeben. Dann gilt
\[
 \bigcup_{i\in I}f_i\text{ ist injektiv}
 \quad\Longleftrightarrow\quad
 \text{die Mengen }f_i[A_i]\text{ sind paarweise disjunkt}.
\]
Insbesondere ist die zusammengefügte Funktion injektiv,
wenn bereits die gewählten Zielmengen $B_i\subseteq B$
paarweise disjunkt sind und $f_i[A_i]\subseteq B_i$ gilt.
Die leere Familie ist zugelassen.
\end{theorem}

\section{Zusätzliche Argumente und freie Werte}

\begin{definition}[Erweiterung einer Funktion]
Seien $A\subseteq A'$, $f\colon A\to B$ und $F\colon A'\to B$.
Dann heißt $F$ eine \emph{Erweiterung} von $f$, wenn $F|_A=f$.
\end{definition}

\begin{theorem}[Erweiterung um ein neues Argument]
\label{b06guide:thm:onepoint}
Seien $f\colon A\rightarrowtail B$, $a\notin A$ und $b\in B$.
Dann gilt
\[
 f\cup\{(a,b)\}\colon A\cup\{a\}\rightarrowtail B
 \quad\Longleftrightarrow\quad b\notin f[A].
\]
Eine injektive Erweiterung um das neue Argument $a$ existiert
also genau dann, wenn $B\setminus f[A]\ne\varnothing$.
\end{theorem}

\begin{theorem}[Allgemeines Erweiterungskriterium]
\label{b06guide:thm:extension}
Seien $A\subseteq A'$ und $f\colon A\rightarrowtail B$.
Dann sind alle injektiven Erweiterungen $F\colon A'\to B$
genau die Funktionen
\[
 F=f\cup h,\qquad
 h\in\operatorname{Inj}\bigl(A'\setminus A,\ B\setminus f[A]\bigr).
\]
Zu jeder solchen Erweiterung ist der zusätzliche Graph eindeutig:
$h=F|_{A'\setminus A}$.
Insbesondere gilt
\[
 \begin{gathered}
 \text{Es existiert eine injektive Erweiterung von }f\text{ auf }A'\\
 \Longleftrightarrow\
 \operatorname{Inj}\bigl(A'\setminus A,\ B\setminus f[A]\bigr)
        \ne\varnothing.
 \end{gathered}
\]
\end{theorem}